% Latin-script Interslavic source component: Work 45 (Kapferer--Noether paper).
% Audited editorial provenance: ../00_SESSION_INDEPENDENT_STATE_v019.json and ../NORMALIZATION_DECISIONS_v019.jsonl.
\documentclass[11pt]{article}
\usepackage{fontspec}
\IfFontExistsTF{Times New Roman}{\setmainfont{Times New Roman}}{\setmainfont{TeX Gyre Termes}}
\IfFontExistsTF{Arial}{\setsansfont{Arial}}{\setsansfont{TeX Gyre Heros}}
\IfFontExistsTF{Consolas}{\setmonofont{Consolas}}{\setmonofont{Latin Modern Mono}}
\usepackage{amsmath,amssymb,mathtools}
\usepackage{amsthm}
\usepackage{geometry}
\usepackage{microtype}
\geometry{a4paper,margin=2.05cm}
\setlength{\parindent}{1.1em}
\setlength{\parskip}{0pt}
\newcommand{\foreign}[1]{#1}
\theoremstyle{plain}
\newtheorem{satz}{Теорем}
\newtheorem{zusatz}{Додаток}
\emergencystretch=220em
\allowdisplaybreaks
\sloppy
\hbadness=10000
\vbadness=10000
\hfuzz=5pt
\vfuzz=12pt
\raggedbottom
\begin{document}

\newpage
\part*{Чест III\hspace{1em}Статја Капферера --- Noether}
\addcontentsline{toc}{part}{Чест III: Статја Капферера --- Noether}

\begin{center}
\Large\textbf{Потрєбне и достаточне условја многократности}\\[0.3em]
\Large\textbf{к Нетеровому фундаменталному теорему}\\[0.5em]
\large\textbf{алгебраичных функциј}\\[1em]
\normalsize\textit{H. Kapferer}\\[0.5em]
\small\textit{(с додатком, сполечнє с E. Noether)}\\[1em]
\normalsize \foreign{Math.\ Ann.}\ \textbf{97} (1927), С.~559--567
\end{center}

\vspace{1em}

Слєдујуче прєдставја заострење Нетерового фундаменталного
теорема в формє потрєбных и достаточных условиј
многократности\textsuperscript{1)}. Комплекс фактов Нетерового
фундаменталного теорема ја, в опрєњу на идеално-теоретично
понимање теорема\textsuperscript{2)}, але без употрєбы
идеално-теоретичных појмов и теоремов ту и в даљшем, сведу до
слєдујучих двех оддєлных тврджениј\textsuperscript{3)}:

\smallskip
\textbf{Теорем Иа.} Дла всакој обчеј нуловеј точкы $x = a$,
$y = b$ двех взајемно простых полиномов $\varphi$, $\psi$ в $x$, $y$,
с либовољными комплексными коефициентами, егзистује принајмење цєло
рационално позитивно число $\varrho$ тако, же модул
$(\varphi, \psi, \mathfrak{p}^{\varrho})$ означа точно ту саму
целост полиномов како модул $(\varphi, \psi, \mathfrak{p}^{\sigma})$,
кде $\mathfrak{p} = (x-a, y-b)$ и $\sigma$ означа либовољно цєло
рационално число вечше од $\varrho$\textsuperscript{4)}.

\smallskip
\textbf{Теорем Иб.} Ако се та појмовно једнозначно опрєдєљена
целост полиномов означи с $\mathfrak{q}$, озирома с
$\mathfrak{q}_{\nu}$, когда иде о $\nu$-тој из $s$ различных нулевых
точек, т.ј. о $x = a_{\nu}$, $y = b_{\nu}$, тогда имаје мєсто:

Потрєбно и достаточно за то, же даны полином $K$ в $x$, $y$ имаје
својство дєлимосци $K \equiv 0(\varphi, \psi)$, јест симултанно
изполњење слєдујучих $s$ конгруенциј
\[
K \equiv 0(\mathfrak{q}_{1}), \quad K \equiv 0(\mathfrak{q}_{2}), \quad\ldots,\quad
K \equiv 0(\mathfrak{q}_{s}).
\]

Хисторично, сравни увод основној расправы Макса Нетера о тој
темє\textsuperscript{5)}, фундаменталны теорем вырастл из покусов
рєшити слєдујучу фундаменталну задачу: при всаком даном полиному
$K$ в $x$, $y$ трєба умєти установити, је-ли он дєлимы с
$(\varphi, \psi)$ или не, т.ј. такоже, егзистујета-ли два друге
полиномы $A$ и $B$ в $x$, $y$ тако, же једначење
\[
K = A\varphi + B\psi
\]
стане идентитето в $x$, $y$.

Та задача не рєшаје се Нетеровым теоремом; послєдњи трєба боље
разумєти само како, хоти фундаменталну, редукцију\textsuperscript{6)}
поставјеној задачи к простєјши, именујуче к пытању о потрєбных и
достаточных условјах за обстојање појединых конгруенциј
$K \equiv 0(\mathfrak{q})$. Когда јест $K \equiv 0(\mathfrak{q})$?
Само когда то послєдње пытање јест одговорјено, јест такоже
изполњена почеткова формулација M.\ Нетера. То јест циљ тој
расправы, кторы такоже фактично јест досєгнуты, и то показовањем
конечно многых условиј многократности, практично без тежкости
контроловатељных, кторе за $K \equiv 0(\mathfrak{q})$ сут
потрєбне и достаточне. Достаточно условје за
$K \equiv 0(\mathfrak{q})$, кторе але не јест такоже потрєбно
условје, јест познано и често употрєбјано и цитовано в литератури,
напримєр в формє:

\smallskip
\textit{За обстојање конгруенције $K \equiv 0(\varphi, \psi)$
достаточно јест, же всака точка прєсека $\varphi = 0$, $\psi = 0$
в полиному $K$ појавја се с ``достаточно'' высокоју многократностју.}

\smallskip
Тај теорем содрживаје се в теорему И како специалны случај; зато же
$\varrho$-кратна точка в $K$ јест равнозначна с
$K \equiv 0(\mathfrak{p}^{\varrho})$; тогда а фортиори јест такоже
$K \equiv 0(\varphi, \psi, \mathfrak{p}^{\varrho})$, т.ј. такоже
$K \equiv 0(\mathfrak{q})$. Але обратно тврджење не имаје мєсто.

\vspace{0.5em}
\noindent\rule{\textwidth}{0.4pt}
\vspace{0.3em}

{\footnotesize
\textsuperscript{1)} За подробно прєдставјење одказују на своју работу с тым самым
насловом в специалном свазку справ Хеиделбержској академије 1927
``\foreign{Beiträge zur Algebra}'', започетом \foreign{A.\ Loewy}, Фреибург и.\ Бр., кторы
содрживаје работы разных авторов.

\textsuperscript{2)} Сравни напр. \foreign{B.\ L.\ van der Waerden}, ``\foreign{Zur Nullstellentheorie der Polynomideale}''.
\foreign{Math.\ Annalen} 96 (1926), С.~203.

\textsuperscript{3)} В \S~1 цитованој расправы ја доказал оба тврджења поједино и директно,
т.ј. без идеално-теоретичных појмов.

\textsuperscript{4)} Буквы $\mathfrak{p}$ и $\mathfrak{q}$ сут выбране зато, абы напоминати
о смысловом согласју означеного с појмами просты идеал
$\mathfrak{p}$ и примарны идеал $\mathfrak{q}$, обычными в новєјшој литератури.

\textsuperscript{5)} \foreign{Max Noether}, ``\foreign{Über einen Satz aus der Theorie der algebraischen Funktionen}'',
\foreign{Math.\ Annalen} 6 (1872).

\textsuperscript{6)} Названа редукција јест приближно сравнима с тој, коју всак зна из
елементарној теорије чисел: својство дєлимосци природного числа
$A$ с таким самым $B$ јест еквивалентно својству дєлимосци $A$ с
$p_{1}^{e_{1}}$, с $p_{2}^{e_{2}}$, \ldots, с $p_{s}^{e_{s}}$, кде
$p_{1}, p_{2}, \ldots, p_{s}$ означајут различне просте дєлитељи
$B = p_{1}^{e_{1}}p_{2}^{e_{2}} \cdots p_{s}^{e_{s}}$. Како аналог
степена простого числа трєба разумєти точно појем ``примарны идеал'',
в тексту означены с $\mathfrak{q}$. [Ср. Е.\ Noether,
``\foreign{Idealtheorie in Ringbereichen}'', \foreign{Math.\ Annalen} 83 (1921), увод.]
}

\newpage
\section*{Главны теорем}

Обєчано рєшење проблема можно кратко карактеризовати тако:
појединому $\mathfrak{q}$ можно придружити систем од $t$ позитивных
цєлых рационалных чисел $(g_{1}), (g_{2}), \ldots, (g_{t})$, а
модулу $(K, \mathfrak{q})$ равно толико цєлых рационалных позитивных
или негативных чисел $(K_{1}), (K_{2}), \ldots, (K_{t})$, тако же
имаје мєсто слєдујучи \textbf{главны теорем}:

\begin{satz}[Главны теорем]\label{satz:kap-II}
Потрєбно и достаточно за $K \equiv 0(\mathfrak{q})$ јест
симултанно изполњење $t$ условиј
\[
(K_{1}) \geq (g_{1}), \quad (K_{2}) \geq (g_{2}), \quad\ldots,\quad (K_{t}) \geq (g_{t}).
\]
При том
\[
(g_{1}) + (g_{2}) + \cdots + (g_{t}) = \mu,
\]
кде $\mu$ означа многократност прєсека, одговарајучу
$\mathfrak{q}$ и опрєдєљену резултантоју.
\end{satz}

В доказу, без ограничења обче валидности, прєдположимо $x = 0$,
$y = 0$ како нулеву точку пара $(\varphi, \psi)$, одговарајучу
$\mathfrak{q}$. Даље, ако $A(x,y)$ јест како-либо полином\textsuperscript{7)}
в $x$, $y$, симбол $(A)$ буде означати число, кторе каже, колико
крат $y$ выступа како фактор в $A(0,y)$; ако але $A(0,y)$
идентично изчезне, тогда $(A)$ ставјаје се равно $\infty$. Числа
$(g_{1}), (g_{2}), \ldots, (g_{t})$, придружене к $\mathfrak{q}$,
опрєдєљајут се слєдујучим системом једначениј, кторы твори
методичну основу поступка:
\begin{equation}\label{eq:kap-1}
\begin{aligned}
f_{1}\cdot g_{1}(0,y):y^{(g_{1})} - g_{1}\cdot f_{1}(0,y):y^{(g_{1})} &= x\cdot h_{1},\\
f_{2}\cdot g_{2}(0,y):y^{(g_{2})} - g_{2}\cdot f_{2}(0,y):y^{(g_{2})} &= x\cdot h_{2},\\
&\vdots \qquad\qquad\vdots\qquad\qquad\vdots\\
f_{r}\cdot g_{r}(0,y):y^{(g_{r})} - g_{r}\cdot f_{r}(0,y):y^{(g_{r})} &= x\cdot h_{r}.
\end{aligned}
\end{equation}
On будује се само из даного пара полиномов $\varphi$, $\psi$:
полиномы $f_{1}, g_{1}$ најпрво опрєдєљајут се како идентичне с
$\varphi, \psi$, евентуално по прєстављању пореда; поредок
должен быти так, же $(f_{1}) \geq (g_{1})$. При тој фиксацији, и
само тогда, $h_{1}$ буде цєло-рационалны.

Полиномы $f_{2}, g_{2}$ опрєдєљајут се како идентичне с
$h_{1}, g_{1}$ при побочном условју $(f_{2}) \geq (g_{2})$. Обче пар
полиномов $f_{r}, g_{r}$ опрєдєља се како идентичны с
$h_{r-1}, g_{r-1}$ при побочном условју $(f_{r}) \geq (g_{r})$.

Систем једначениј \eqref{eq:kap-1} можно истинно безкончно
продолжати в објасњеном способу; главны теорем але потрєбује
створити само прве $t$ једначења система, кде $t$ јест опрєдєљено
слєдујучим фактом:

\smallskip
\textbf{Помочны теорем I.} Егзистује природно число $t$ [и то
$t \leq \mu$, под $\mu$ разумєваје се многократност прєсека, в којој
точка $x = 0$, $y = 0$ выступа в $\varphi = 0$, $\psi = 0$] тако, же
прве $t$ числа, опрєдєљене с \eqref{eq:kap-1}, именујуче
$(g_{1}), (g_{2}), \ldots, (g_{t})$, сут всє вечше од нуле, покым
једночасно $(g_{t+1})$ и всє слєдујуче сут равне нули.

С урєдњењем продуктного теорема резултантов, многократност
прєсека в точце $x = 0$, $y = 0$ за $f_{r} = 0$, $g_{r} = 0$ буде
равна
\[
\mu - (g_{1}) - (g_{2}) - \cdots - (g_{r}),
\]
из чего слєдује прєрыв по највыше $\mu$ кроках. Из тога једночасно
слєдује факт
\begin{equation}\label{eq:kap-2}
\mu = (g_{1}) + (g_{2}) + \cdots + (g_{t}),
\end{equation}
формула значителна зато, же она обче даје рационалну и практично
легко извршиму методу опрєдєљења многократности прєсека, когда
одговарајуча нулева точка јест познана.

Когда числа $(g_{i})$ сут тако познане, числа $(K_{i})$ опрєдєљајут
се слєдујучим системом једначениј, одговарајучим к \eqref{eq:kap-1}:
\begin{equation}\label{eq:kap-3}
\begin{aligned}
K_{1}\cdot g_{1}(0,y):y^{(g_{1})} - g_{1}\cdot K_{1}(0,y):y^{(g_{1})} &= x\cdot K_{2},\\
K_{2}\cdot g_{2}(0,y):y^{(g_{2})} - g_{2}\cdot K_{2}(0,y):y^{(g_{2})} &= x\cdot K_{3},\\
&\vdots \qquad\qquad\vdots\qquad\qquad\vdots\\
K_{r}\cdot g_{r}(0,y):y^{(g_{r})} - g_{r}\cdot K_{r}(0,y):y^{(g_{r})} &= x\cdot K_{r+1}.
\end{aligned}
\end{equation}
$K_{1}$ опрєдєља се како идентичны с $K$. Јасно јест, же $K_{2}$
буде цєло-рационалны тогда и само тогда, когда
$(K_{1}) \geq (g_{1})$. Ако то послєдње јест изполњено, даље
слєдује: $K_{3}$ буде цєло-рационалны тогда и само тогда, когда
$(K_{2}) \geq (g_{2})$. Обче, ако $K_{r}$ уж јест цєло-рационалны,
$K_{r+1}$ буде цєло-рационалны тогда и само тогда, когда
$(K_{r}) \geq (g_{r})$.

Доказ главного теорема (теорема~\ref{satz:kap-II}) тераз опрє се на
даљши

\smallskip
\textbf{Помочны теорем II.} Ако се
$(f_{i}, g_{i}, \mathfrak{p}^{e})$ постави равно
$\mathfrak{q}^{(i)}$, тогда за всако $i$ из цєлочиселној серије
$1, 2, \ldots$ \textit{ин инф.} имаје мєсто теорем:

За конгруенцију
\[
K_{i} \equiv 0(\mathfrak{q}^{(i)})
\]
потрєбно и достаточно јест симултанно изполњење двех условиј:
\[
(K_{i}) \geq (g_{i}) \qquad \text{и} \qquad K_{i+1} \equiv 0(\mathfrak{q}^{(i+1)}).
\]

Доказ\textsuperscript{8)} помочнохо теорема добыва се комбиновањем
системов једначениј \eqref{eq:kap-1} и \eqref{eq:kap-3}. Ако се
урєдни, же $g_{t+1}$ по помочном теорему И в нулевој точце вече не
изчезне, т.ј. же модул $\mathfrak{q}^{(t+1)}$ состоји се из всих
полиномов (јест јединичны идеал), и зато
$K_{t+1} \equiv 0(\mathfrak{q}^{(t+1)})$ стане тривиално, тогда
$t$-кратно употрєбјење помочнохо теорема II даје доказ главного
теорема.

\smallskip
Без доказов, кторе находет се в мојеј горе цитованој расправи,
додају јешче два додатка к главному теорему:

\begin{zusatz}[I]
Число $t$ главного теорема јест идентично с принајмење\textsuperscript{9)}
позитивным числом $\nu$ такохо рода, же
$x^{\nu} \equiv 0(\mathfrak{q})$.
\end{zusatz}

\begin{zusatz}[II]
Системы једначениј \eqref{eq:kap-1} и \eqref{eq:kap-3}, и те, кторе
из њих взникајут прєстављањем $x$ с $y$, при входној комбинацији
давајут рационалну методу точного числовного опрєдєљења
``належечего'' к $\mathfrak{q}$ експонента
$\varrho$, т.ј. того минималного експонента $\varrho$ из теорема Иа.
\end{zusatz}

Е.\ Бертини\textsuperscript{10)} дал обчу формулу за $\varrho$.
Же та формула але не всегда даје минималну врєдност $\varrho$,
показано јест в цитованој расправи на прикладу
$\varphi = x^{3} + y^{4}$, $\psi = x^{2} + y^{3}$. Наконец трєба
јешче замєтити, же и сама формула Бертинија може быти по новому
основана нашим главным теоремом.

\vspace{0.5em}
\noindent\rule{\textwidth}{0.4pt}
\vspace{0.3em}

{\footnotesize
\textsuperscript{7)} Ако за $A(x,y)$ допустимо такоже дробно-рационалне
функције $C(x,y):D(x,y)$, тогда $(A) = (C) - (D)$; в послєдњем
случају може $(A)$ имєти и негативне врєдности.

\textsuperscript{8)} Ср. такоже доказ в замєтце И ``Додатка''.

\textsuperscript{9)} Ту се зато добывајут точне експоненты в противности к
оцєнкам господжице Grete Hermann [``\foreign{Die Frage der endlich vielen Schritte
in der Theorie der Polynomideale}'' (с употрєбоју посмртно оставјеных
теоремов К.\ Хентзелт), \foreign{Math.\ Annalen} 95 (1926)], кде при
много обчејшој поставјености проблема давајут се само горне границе.

\textsuperscript{10)} Е.\ Бертини, \foreign{Math.\ Annalen} 34 (1889).
}

\newpage
\section*{Теорем III и додаток}

Интересна специализација главного теорема слєдује из ограничајучего
прєдположења, же нулева точка одговарајучего
$\mathfrak{q} = (\varphi, \psi, \mathfrak{p}^{e})$ јест проста
точка принајмење в једној из двех кривуљ $\varphi = 0$, $\psi = 0$.
Ако она јест напр. проста точка в $\psi$, тогда можно показати, же
число $t$ главного теорема јест точно равно $\mu$ и же $t$
неравности главного теорема можно замєнити једино условјем
многократности. То води к слєдујучему обчему теорему:

\begin{satz}[Теорем III]\label{satz:kap-III}
Ако всє точкы прєсека $\varphi = 0$, $\psi = 0$ сут просте точкы в
$\psi$, тогда, прєдположивши же $K$ јест взајемно просты с $\psi$,
за обстојање конгруенције $K \equiv 0(\varphi, \psi)$ потрєбно и
достаточно јест, же $K$ изчезаје в всих тых точках, але тако, же
многократност прєсека в пару кривуљ $K = 0$, $\psi = 0$ в всакој
из њих јест принајмење така велика како в пару кривуљ
$\varphi = 0$, $\psi = 0$.
\end{satz}

Тај теорем III јест јешче из особној причины значителны: условје
многократности, кторе он потрєбује, јест тако, же јего изполњење
или неизполњење може быти установјено конечно-кратным
диференцовањем, независно од формулы \eqref{eq:kap-2}, опрєтој на
\eqref{eq:kap-1}. То јест непосрєдње јасно, ако цитују слєдујучи
``теорем'', кторы ја поставил в 1923 року\textsuperscript{11)}:

\smallskip
\textit{Ако $A$ и $B$ сут како-либо два взајемно просте полиномы в
$x$, $y$, а $P$ јест обча точка, ктора јест проста точка од $B$,
тогда многократност тој точкы како точкы прєсека јест тогда и само
тогда равна числу $\mu$, когда многократност тој самеј точкы в пару
кривуљ $A' = 0$, $B = 0$ јест точно $\mu - 1$; при том означа}
\[
A' = \frac{\partial A}{\partial x} \cdot \frac{\partial B}{\partial y} -
\frac{\partial A}{\partial y} \cdot \frac{\partial B}{\partial x}.
\]

Наконец трєба јешче споменути, же главны теорем
(теорем~\ref{satz:kap-II}) с малоју модификацијоју може быти
прєнесены такоже на хомогене полиномы в 3 вариаблах $x$, $y$, $z$.
То прєнесење удаје се особито елегантно в горе наведном специалном
случају. Именно можно показати:

\begin{zusatz}[к теорему III]
Теорем~\ref{satz:kap-III} оставаје цєлком незмєњены, и по словах,
ако под $\varphi$, $\psi$, $K$ разумєвајут се хомогене полиномы в
трех вариаблах\textsuperscript{12)}.
\end{zusatz}

\vspace{0.5em}
\noindent\rule{\textwidth}{0.4pt}
\vspace{0.3em}

{\footnotesize
\textsuperscript{11)} Х.\ Капферер, ``\foreign{Über die Multiplizität der Schnittpunkte von zwei algebraischen
Kurven}''. \foreign{Jahresber.\ der D.\ M.-V.} 32 (1923).

\textsuperscript{12)} Пытање о условјах за обстојање хомогене конгруенције в
трех вариаблах $K(x,y,z) \equiv 0(\varphi(x,y,z), \psi(x,y,z))$ јест,
хоти под силно специализоваными прєдположењами, такоже јадро
изслєдовања \foreign{A.\ Ostrowski}, ``\foreign{Über die Maxwellsche Erzeugung der
Kugelfunktionen}'', \foreign{Jahresb.\ der D.M.-V.} 33 (1925). В тој
нотє прєдположено јест: $\varphi$ разпадаје се на саме линеарне
факторы; $\varphi$ и $K$ сут једнакохо пореда; $\psi = x^{2} + y^{2} + z^{2}$.
Зато же ту $\psi$ јест цєлком без сингуларностиј, главны теорем
Островскијевој ноты можно разумєти како специалны случај нашего
теорема III.
}

\newpage
\section*{Додаток, сполечнє с E. Noether\textsuperscript{13)}}

Систем једначениј \eqref{eq:kap-1} даје једночасно рационалну
методу за опрєдєљење полного система остатков по $\mathfrak{q}$,
под чим трєба разумєти полны систем прєдставитељев конечно многых,
в односу к области коефициентов $P$, модуло $\mathfrak{q}$ линеарно
независных остатковых класов. С тым поједине крокы в доказу
главного теорема ставајут много прозрєчне; једночасно како
побочны резултат выходи факт, же резултанта-многократност $\mu$
согласује се с должиноју идеала $\mathfrak{q}$;

\vspace{0.5em}
\noindent\rule{\textwidth}{0.4pt}
\vspace{0.3em}

{\footnotesize
\textsuperscript{13)} Идеално-теоретично толковање походит од Е.\ Noether, а основны
доказ (5) и (6) од Х.\ Капферер.

\textsuperscript{14)} Поље $P$ трєба, без ограничења обче валидности, прєдположити
алгебраично затвореным; в том обсегу фактично имајут мєсто всє попрєдне
теоремы, с изкључењем диференциацијского условја. Должина
$\mathfrak{q}$ јест такоже опрєдєљена како должина композицијној
серије идеалов, идуче од $\mathfrak{p}$ до $\mathfrak{q}$. Ср.
Е.\ Noether, \foreign{Jahresber.\ d.\ d.\ Math.-Ver.} 34 (1925), С.~101
(коса пагинација); подробнєје у Х.\ Грелл, \foreign{Math.\ Annalen} 97
(1927), С.~529.

Доказы в литератури за согласје обојих чисел многократности (при
$n$ вариаблах) сут зело сложне; најпростєјши јест јешче непубликованы
посмртно оставјены доказ Хентзелта.
}

\vspace{1em}
т.ј. с рангом прстена остатковых класов или с числом в односу к $P$
линеарно независных остатковых класов\textsuperscript{14)}.

Полны систем остатков јест, с означењами из \eqref{eq:kap-1} и
помочнохо теорема II, карактеризованы слєдујучим

\begin{satz}[Теорем IV]\label{satz:kap-IV}
Полны систем остатков по $\mathfrak{q}$ даны јест полиномами
\[
h_{t},\ h_{t}y,\ldots,h_{t}y^{(g_{t})-1};\ \ldots;\
h_{i},\ h_{i}y,\ldots,h_{i}y^{(g_{i})-1};\ \ldots;\
h_{1},\ h_{1}y,\ldots,h_{1}y^{(g_{1})-1}.
\]
При том всакы раз
$h_{i}, h_{i}y, \ldots, h_{i}y^{(g_{i})-1}$ даје полны
систем остатков од $\mathfrak{q}^{(i+1)}$ до
$\mathfrak{q}^{(i)}$.
\end{satz}

Из дефиниције слєдује непосрєдње:
$\mathfrak{q}^{(i+1)} = (\mathfrak{q}^{(i)}, h_{i})$, зато
$\mathfrak{q}^{(i)} \equiv 0(\mathfrak{q}^{(i+1)})$.
Теорем~\ref{satz:kap-IV} буде доказаны, ако показано јест:
\begin{equation}\label{eq:kap-4}
h_{i}x \equiv 0(\mathfrak{q}^{(i)}); \qquad
h_{i}y^{(g_{i})} \equiv 0(\mathfrak{q}^{(i)})
\end{equation}
и
\begin{equation}\label{eq:kap-5}
\text{из } h_{i}F(y) \equiv 0(\mathfrak{q}^{(i)})
\quad \text{слєдује} \quad F(y) \equiv 0(y^{(g_{i})}).
\end{equation}

Релације \eqref{eq:kap-4} именно говорет, же линеарне комбинације
$h_{i}, h_{i}y, \ldots, h_{i}y^{(g_{i})-1}$ изчерпавајут
остаткове класы од $\mathfrak{q}^{(i+1)}$ до
$\mathfrak{q}^{(i)}$, покым \eqref{eq:kap-5} выражаје линеарну
независност. Зато же $\mathfrak{q}^{(t+1)}$ ставаје јединичным
идеалом; то тераз слєдује из конечној должины $\mathfrak{q}$,
из тога добыва се доказ теорема~\ref{satz:kap-IV} и једночасно
факт, же должина $\mathfrak{q}$ ставаје равна
$(g_{1}) + \cdots + (g_{t})$, зато по \eqref{eq:kap-2} согласује се
с резултанта-многократностју.

Прва из релациј \eqref{eq:kap-4} чита се из \eqref{eq:kap-1};
друга добыва се слєдујучим способом:

Имамо $h_{i}g_{i} \equiv 0(\mathfrak{q}^{(i)})$; зато заради прве
релације \eqref{eq:kap-4} такоже
$h_{i}g_{i}(0,y) \equiv 0(\mathfrak{q}^{(i)})$, или
$h_{i}y^{(g_{i})}(g_{i}(0,y):y^{(g_{i})}) \equiv 0(\mathfrak{q}^{(i)})$,
из чего $h_{i}y^{(g_{i})} \equiv 0(\mathfrak{q}^{(i)})$
слєдује заради $g_{i}(0,y):y^{(g_{i})} \not\equiv 0(y)$.

Пред доказом \eqref{eq:kap-5} дајмо опазку непосрєдње читану из
\eqref{eq:kap-1}: $f_{i}$ и $g_{i}$ могут имєти како обчи дєлитељ
само полином $d_{i}(y)$ в $y$, не дєлимы с $y$:
$f_{i} = d_{i}(y)\bar{f}_{i}$, $g_{i} = d_{i}(y)\bar{g}_{i}$, и
$\bar{f}_{i}$ и $\bar{g}_{i}$ сут взајемно просте. Заради
$d_{i}(y) \not\equiv 0(\mathfrak{p})$ буде $\mathfrak{q}^{(i)}$
такоже равно $(\bar{f}_{i}, \bar{g}_{i}, \mathfrak{p}^{e})$;
при том то буде равно
$(\bar{f}_{i}, \bar{g}_{i}, \mathfrak{p}^{\sigma})$ за всако
$\sigma \geq \varrho$, како дєлитељ $\mathfrak{q}$ за всако
тако $\sigma$; зато $\mathfrak{q}^{(i)}$ јест примарна компонента
$(\bar{f}_{i}, \bar{g}_{i})$, належеча к нулевој точце.

Нехај $\gamma_{1}, \delta_{1}; \ldots; \gamma_{r}, \delta_{r}$ сут
нулеве точкы $(\bar{f}_{i}, \bar{g}_{i})$, различне од
$x = 0$, $y = 0$, и нехај $\tau_{1}, \ldots, \tau_{r}$ сут експоненты
одговарајучих примарных компонентов $(\bar{f}_{i}, \bar{g}_{i})$;
даље в продукту
\[
P = ((x-\gamma_{1}) + u(y-\delta_{1}))^{\tau_{1}} \cdots
((x-\gamma_{r}) + u(y-\delta_{r}))^{\tau_{r}}
\]
неизвєстну $u$ выберемо како елемент из $P$ тако, же остане
$P \not\equiv 0(\mathfrak{p})$. Ако се $G(x,y)$ постави равно
продукту специализованого $P$ с $d_{i}(y)$, тогда
$G(x,y) \not\equiv 0(\mathfrak{p})$ и
$\mathfrak{q}^{(i)} \cdot G(x,y) \equiv 0(f_{i}, g_{i})$.

Из прєдположења $h_{i}F(y) \equiv 0(\mathfrak{q}^{(i)})$
слєдује
$h_{i}F(y)G(x,y) \equiv 0(f_{i}, g_{i})$, или, сполочно
с \eqref{eq:kap-1}, в формє идентитетов:
\[
h_{i}F(y)G = Af_{i} + Bg_{i}; \qquad
h_{i}x = (g_{i}(0,y):y^{(g_{i})})\cdot f_{i}
- (f_{i}(0,y):y^{(g_{i})})g_{i}.
\]
То даје але
\[
\bar{f}_{i}(xA - F(y)H) \equiv 0(\bar{g}_{i}); \qquad
H = G(x,y) \cdot g_{i}(0,y) : y^{(g_{i})} \not\equiv 0(\mathfrak{p}).
\]
Из взајемној простости $\bar{f}_{i}$ и $\bar{g}_{i}$ слєдује даље
\[
xA - F(y)H \equiv 0(\bar{g}_{i});
\]
тако $F(y)H \equiv 0(\bar{g}_{i}, x)
\equiv 0(\bar{g}_{i}, x, \mathfrak{p}^{\sigma})
= (y^{(g_{i})}, x)$, и с тым $F(y) \equiv 0(y^{(g_{i})})$.
Тым јест \eqref{eq:kap-5}, и зато теорем~\ref{satz:kap-IV}, доказаны.

\smallskip
\textbf{Замєтка I.} Релације \eqref{eq:kap-4} и \eqref{eq:kap-5}
сполочно с помочным теоремом II можно записати како взајемно
реципрочне релације мед квоциентами идеалов\textsuperscript{15)};
при том доказ другој релације даје једночасно доказ помочнохо
теорема II:
\begin{equation}\label{eq:kap-6}
\mathfrak{q}^{(i)} : \mathfrak{q}^{(i+1)}
= (\mathfrak{q}^{(i)}, x) = (y^{(g_{i})}, x); \qquad
\mathfrak{q}^{(i)} : (\mathfrak{q}^{(i)}, x)
= \mathfrak{q}^{(i+1)}.\qquad\text{\textsuperscript{16)}}
\end{equation}
Прва релација прєдставја сјединење \eqref{eq:kap-4} и
\eqref{eq:kap-5}; к другој трєба замєтити: в \eqref{eq:kap-4}
было доказано
$x \cdot \mathfrak{q}^{(i+1)} \equiv 0(\mathfrak{q}^{(i)})$;
обратно добыва се аналогично како \eqref{eq:kap-5}. Из
$xM \equiv 0(\mathfrak{q}^{(i)})$ слєдује како выше:
\[
xMG \equiv 0(f_{i}, g_{i}); \quad \text{тако }\quad
xMG \cdot g_{i}(0,y) : y^{(g_{i})}
\equiv 0(xh_{i}, g_{i}).
\]
Зато же $g_{i} \not\equiv 0(x)$ по дефиницији, именујуче
$(f_{i}) \geq (g_{i})$, из тога выходи
\[
M \cdot (G \cdot g_{i}(0,y) : y^{(g_{i})})
\equiv 0(h_{i}, g_{i}) \equiv 0(\mathfrak{q}^{(i+1)}),
\quad\text{и зато}\quad M \equiv 0(\mathfrak{q}^{(i+1)}).
\]

Же условје помочнохо теорема II јест потрєбно, слєдује тераз тако:
$K_{i} \equiv 0(\mathfrak{q}^{(i)})$ даје
$K_{i}(0,y) \equiv 0(y^{(g_{i})}, x)$; зато
$(K_{i}) \geq (g_{i})$ и с тым, по \eqref{eq:kap-3},
\[
x \cdot K_{i+1} \equiv 0(\mathfrak{q}^{(i)});
\]
зато по другој релацији \eqref{eq:kap-6} буде
$K_{i+1} \equiv 0(\mathfrak{q}^{(i+1)})$.

Же условје јест достаточно, чита се непосрєдње; конечно-кратно
употрєбјење помочнохо теорема II даје але главны теорем. Такоже
итерација другој релације \eqref{eq:kap-6} даје додаток И, именујуче
$\mathfrak{q} : x^{t} = \mathfrak{q}^{(t+1)}$, зато јединичны
идеал, и $t$ јест најнижи так степен, заради
$\mathfrak{q} : x^{t-1} = \mathfrak{q}^{(t)}$.

\textbf{Замєтка II.} За полином $K$, не дєлимы с $\mathfrak{q}$,
можно јешче тврдити: ако
$K \equiv 0(\mathfrak{q}^{(i+1)})$,
$\not\equiv 0(\mathfrak{q}^{(i)})$, тогда
$(K, \mathfrak{q}^{(i)}) = (h_{i}y^{(K)}, \mathfrak{q}^{(i)})$.
Зато же по теорему~\ref{satz:kap-IV}
$K \equiv h_{i}y^{\lambda}c(y)$ модуло
$(\mathfrak{q}^{(i)})$ с $c(y) \not\equiv 0(y)$; тако выходи
\[
(K, \mathfrak{q}^{(i)}) = (h_{i}y^{\lambda}, \mathfrak{q}^{(i)}).
\]
При том $y^{(g_{i})-\lambda} \cdot K \equiv 0(\mathfrak{q}^{(i)})$,
и то по \eqref{eq:kap-5} јест најнижи тако степен. По помочном
теорему II јест але $(g_{i}) - (K)$ тој најнижи степен; зато
$\lambda$ ставаје равно $(K)$.

\vfill
\begin{flushright}
(Пришло 24.\ 10.\ 1926.)
\end{flushright}

\vspace{0.5em}
\noindent\rule{\textwidth}{0.4pt}
\vspace{0.3em}

{\footnotesize
\textsuperscript{15)} Под квоциентом $\mathfrak{a}:\mathfrak{b}$, како познано,
разумєје се целост полиномов $c(x,y)$ такых, же
$\mathfrak{b}c \equiv 0(\mathfrak{a})$.

\textsuperscript{16)} Из тых двех взајемно реципрочных релациј једна условјује
другу по обчеј теорији иредуцибилных идеалов; ср. Мацаулаы,
\textit{\foreign{Modular Systems}}, Нр.\,72 и 73. \foreign{Cambridge Tracts} 19 (1916).
}

\end{document}
